\section{Separation of subgroup supports}\label{sec:support}

\Cref{cor:finite-gluing} bounds the interaction between a subgroup and a
normal quotient. We now use those bounds to control specialization for every
embedded subgroup. We follow the finite-localization and support-splitting
strategy of Balmer--Sanders \cite[proof of Theorem~9.1, \textup{(B)}
$\Rightarrow$\textup{(A)}]{BS}: a specialization violating the desired bound
would disconnect the irreducible support of a compact object. Here the
cyclic-length estimate isolates an individual subgroup component by
controlling its interactions with normal conjugates above and below it.

\subsection{Subgroup layers and their right adjoints}

Fix a finite \(p\)-group \(G\) and an integer \(m\ge0\). Put
\[
 E=L_m^fS,\qquad e=\operatorname{inf}_G E,
 \qquad \mathcal C=\mathcal C_m^G.
\]
For \(J\le G\), write \(N_J=N_G(J)\), \(W_J=N_J/J\), and
\[
 \mathcal D_J=\Fun(BW_J,\Sp_{L_m^f}).
\]
Let \(\phi_J:\mathcal C\to\mathcal D_J\) be geometric \(J\)-fixed points,
retaining the underlying coherent \(W_J\)-action. We use one representative
of each conjugacy class when indexing these functors. The notation
\(J\sim K\) means conjugacy in \(G\).

\begin{lemma}[Layer adjunctions]\label{lem:layer-adjunction}
The functor \(\phi_J\) is strong symmetric monoidal and has a fully faithful
right adjoint \(F_J\). Write
\(\mathcal A_J=F_J(\one_{\mathcal D_J})\). Then:
\begin{enumerate}[label=\textup{(\roman*)}]
\item The underlying geometric fixed points of its unit value are
\begin{equation}\label{eq:layer-unit}
 \Phi^K \mathcal A_J\simeq
 \bigoplus_{\substack{J'\sim J\\J'\trianglelefteq K}}
 U_{K/J'}(E).
\end{equation}
For arbitrary \(M\in\mathcal D_J\), choose, for each normal conjugate
\(J'=gJg^{-1}\trianglelefteq K\), the diagram \(M_g\) obtained by
restricting \(M\) along
\[
 K/J'\longrightarrow W_J,\qquad xJ'\longmapsto g^{-1}xgJ.
\]
The formula for \(\Phi^K F_JM\) has the same index set, with summand
\(\Phi^{K/J'}\Bor_{K/J'}(M_g)\). In particular,
\(\phi_KF_JM=0\) if no conjugate of \(J\) is normal in \(K\).
\item For every dualizable \(z\in\mathcal C\) and every
\(M\in\mathcal D_J\), the projection map is an equivalence
\begin{equation}\label{eq:layer-projection}
 z\otimes F_JM\simeq F_J(\phi_Jz\otimes M).
\end{equation}
\item The object \(\mathcal A_J\) is a commutative algebra, and \(F_JM\) is a
unital \(\mathcal A_J\)-module.
\end{enumerate}
\end{lemma}

\begin{proof}
First work before chromatic localization. With \(N=N_J\), \(W=W_J\),
factor the functor as
\[
 \Sp^G\xrightarrow{\operatorname{Res}_N^G}\Sp^N
 \xrightarrow{\Phi_N^J}\Sp^W
 \longrightarrow\Fun(BW,\Sp).
\]
Here the middle functor retains genuine quotient symmetry. Let
\(\mathscr F_J=\{L\le N:J\not\le L\}\), and let
\(\widetilde E\mathscr F_J\) be the cofiber of
\((E\mathscr F_J)_+\to S\). Normal geometric localization identifies
\(\Sp^W\) with the full subcategory of genuine \(N\)-spectra whose
geometric fixed points vanish at subgroups not containing \(J\)
\cite[Section~2\textup{(H)}]{BS}. Its fully faithful right adjoint is
\[
 \rho_J(Y)=\widetilde E\mathscr F_J\otimes\pi^*Y,
 \qquad \pi:N\longrightarrow W,
\]
where \(\pi^*\) is genuine inflation. Consequently the required right
adjoint before chromatic localization is the explicit composite
\begin{equation}\label{eq:layer-construction}
 F_JM=\operatorname{Coind}_{N}^{G}\bigl(\rho_J\Bor_W(M)\bigr).
\end{equation}

For completeness, compute its geometric \(K\)-fixed points. Coinduction
equals induction for finite groups. The fixed-orbit formula therefore gives
one term for each coset \(gN\in(G/N)^K\). Such a coset satisfies
\(g^{-1}Kg\le N\). The factor \(\widetilde E\mathscr F_J\) kills this
term unless \(J\le g^{-1}Kg\). Thus the surviving cosets correspond
bijectively to conjugates \(J'=gJg^{-1}\) normal in \(K\). Above such a
subgroup, normal geometric inflation has the quotient fixed-point formula;
restriction of a cofree spectrum is again cofree. Its coefficient diagram
is the restriction \(M_g\) specified in the statement, so the resulting
term is \(\Phi^{K/J'}\Bor_{K/J'}(M_g)\). For the constant diagram
\(E\), this is exactly \(U_{K/J'}(E)\), proving
\eqref{eq:layer-unit}. If the index set is empty, the fixed points vanish
for every \(M\).

All these constructions retain the residual actions. More explicitly,
\(N_G(K)\) acts on the surviving cosets by left translation, or equivalently
on the normal conjugates \(J'\) by conjugation. It transports each term by
the conjugation and restriction maps in \eqref{eq:layer-construction}; the
action factors coherently through \(W_K\). Formula
\eqref{eq:layer-unit} records the underlying spectrum of this diagram.
At \(K=J\), a surviving conjugate \(J'\le J\) has the same order as
\(J\), so \(J'=J\). Only the identity coset of \(N\) remains. The
counit \(\phi_JF_JM\to M\) then identifies with the cofree counit on
underlying \(W_J\)-diagrams and is an equivalence. This proves full
faithfulness, including when \(J\) is not normal in \(G\).

The unlocalized right adjoint already takes \(E\)-local diagrams to
\(\mathcal C\). Indeed, for such an \(M\) and every genuine \(G\)-spectrum
\(X\), adjunction gives
\begin{align*}
 \Map(e\otimes X,F_JM)
 &\simeq\Map(E\otimes\phi_JX,M)\\
 &\simeq\Map(\phi_JX,M)
 \simeq\Map(X,F_JM).
\end{align*}
These equivalences identify precomposition with the localization map
\(X\to e\otimes X\); this is exactly the criterion for \(F_JM\) to be
local. Hence the adjunction and its formulas restrict to the categories
under consideration.

In the restricted adjunction, duality gives, for \(X\in\mathcal C\),
\begin{align*}
 \Map(X,z\otimes F_JM)
 &\simeq\Map(z^\vee\otimes X,F_JM)\\
 &\simeq\Map(\phi_J(z^\vee)\otimes\phi_JX,M)\\
 &\simeq\Map(\phi_JX,\phi_Jz\otimes M).
\end{align*}
This proves \eqref{eq:layer-projection}. It uses dualizability in
\(\mathcal C\), and places no finiteness condition on \(M\).
The projection equivalence also identifies the unit on \(z\otimes w\)
with \(z\) tensored with the unit on \(w\): their adjoints agree by the
triangle identity
\(\varepsilon_{\phi_Jw}\circ\phi_J(\eta_w)=\id_{\phi_Jw}\).
Finally,
the right adjoint of a strong symmetric monoidal functor is lax symmetric
monoidal; its unit value is therefore a commutative algebra and each of its
values is a unital module over that algebra.
\end{proof}

\subsection{A finite splitting argument}

For a dualizable object \(z\) and a subgroup \(B\le G\), the unit
\(z\to F_B\phi_Bz\) becomes an equivalence after applying \(\phi_B\).
To split this map, we will confine its target to that layer and show
that its fiber remains zero there while the other layers are removed.
The two hypotheses below provide these separate controls through the
projection formula.

\begin{lemma}[Splitting an isolated layer]\label{lem:layer-splitting}
Let \(z\in\mathcal C\) be dualizable and let \(B\le G\). Suppose that
\begin{align}
 \phi_K(z\otimes \mathcal A_B)&=0 &&(K\not\sim B),\label{eq:outgoing-orthogonality}\\
 \phi_B(z\otimes \mathcal A_J)&=0 &&(J\not\sim B).\label{eq:incoming-orthogonality}
\end{align}
Then the adjunction unit splits as
\[
 z\simeq X\oplus Q,\qquad Q=F_B\phi_Bz\simeq z\otimes \mathcal A_B,
\]
with \(\phi_BX=0\), \(\phi_KQ=0\) for \(K\not\sim B\), and
\(X\otimes Q=0\). If \(z\) is compact, both summands are compact.
\end{lemma}

\begin{proof}
Let \(X\) be the fiber of \(z\to Q\). By the projection formula,
\[
 X=z\otimes w_0,\qquad
 w_0=\operatorname{fib}(\one_{\mathcal C}\to \mathcal A_B).
\]
Full faithfulness of \(F_B\) gives \(\phi_BX=0\).

List the other subgroup conjugacy classes as \(J_1,\ldots,J_t\), in
nondecreasing order of subgroup size. Set \(X_0=X\) and inductively form
fiber sequences
\begin{equation}\label{eq:layer-stripping}
 X_s\longrightarrow X_{s-1}\longrightarrow
 Y_s:=F_{J_s}\phi_{J_s}X_{s-1}
 \qquad(1\le s\le t).
\end{equation}
We verify that the process preserves three properties:
\(X_s=z\otimes w_s\) for some \(w_s\), its \(B\)-component is zero,
and its \(J_1,\ldots,J_s\)-components are zero.

If \(X_{s-1}=z\otimes w_{s-1}\), the projection formula identifies
\[
 Y_s\simeq z\otimes F_{J_s}\phi_{J_s}w_{s-1}.
\]
It also identifies the map to \(Y_s\) with \(z\) tensored with the
adjunction unit on \(w_{s-1}\). Thus its fiber has the same form.
The \(J_s\)-component of the fiber vanishes by full faithfulness.
A previously removed subgroup has smaller order, or the same order in a
different conjugacy class. It cannot contain a conjugate of \(J_s\), so
\Cref{lem:layer-adjunction}\textup{(i)} shows that its component of
\(Y_s\) is zero. Its vanishing therefore persists in the fiber.

To check the \(B\)-component, put
\(M_s=F_{J_s}\phi_{J_s}w_{s-1}\). This is an \(\mathcal A_{J_s}\)-module.
For any unital \(A\)-module \(M\), the unit and action exhibit
\(M\) as a retract of \(A\otimes M\). Consequently
\(\phi_Bz\otimes\phi_BM_s\) is a retract of
\[
 (\phi_Bz\otimes\phi_B\mathcal A_{J_s})\otimes\phi_BM_s=0
\]
by \eqref{eq:incoming-orthogonality}. Hence \(\phi_BY_s=0\), and
\(\phi_BX_s=0\) follows. This argument applies to the possibly
noncompact objects \(w_s\) and \(Y_s\).

At the end every subgroup component of \(X_t\) is zero; geometric
detection gives \(X_t=0\). On the other hand,
\eqref{eq:outgoing-orthogonality} and adjunction give, for each \(s\),
\[
 \Map(Q,Y_s)
 \simeq\Map(\phi_{J_s}Q,\phi_{J_s}X_{s-1})=0.
\]
Here \(\Map\) denotes the mapping spectrum, so this vanishing holds in
every degree. Applying \(\Map(Q,-)\) to
\eqref{eq:layer-stripping} and descending from \(X_t=0\) gives
\(\Map(Q,X)=0\). In particular the boundary \(Q\to\Sigma X\)
in the original fiber sequence is zero, proving the splitting.

The stated component vanishings follow from construction and
\eqref{eq:outgoing-orthogonality}. They imply
\(\Phi^K(X\otimes Q)=0\) for every \(K\), hence \(X\otimes Q=0\).
When \(z\) is compact, its retracts \(X\) and \(Q\) are compact as well.
\end{proof}

\subsection{Minimal-height separation}

\begin{lemma}[Chromatic orthogonality]\label{lem:layer-orthogonality}
Let \(z\in\mathcal C_m^G\) be the image of a finite genuine
\(G\)-spectrum, and let \(n\le m\). Suppose that \(B\le G\) satisfies
\[
 m-n<\ell(G,B),
\]
and that every \((K,k)\in\supp(z)\) with \(k<m\) satisfies
\[
 \ell(G,K)\le k-n.
\]
Then \(z\) and \(B\) satisfy both orthogonality conditions in
\Cref{lem:layer-splitting}.
\end{lemma}

\begin{proof}
There is no point \((B,k)\in\supp(z)\) with \(k<m\): otherwise the
second hypothesis would contradict the first. Thus only the test at height
\(m\) can be nonzero on \(\Phi^Bz\).

Fix \(K\not\sim B\). In a summand of \(\Phi^K\mathcal A_B\), let
\(B'\sim B\) be normal in \(K\), and put
\(d=\ell(K,B')=\ell(K/B',1)\). The equality follows by the subgroup
correspondence for \(K\to K/B'\). Since \(K\not\sim B\), this quotient
is nontrivial and \(d>0\). Consider a height
\(0\le k\le m\) with \((K,k)\in\supp(z)\). If \(k=m\), then
\(m<k+d\). If \(k<m\), concatenate cyclic subnormal chains and use
conjugacy invariance to obtain
\[
 m-n<\ell(G,B')
 \le\ell(K,B')+\ell(G,K)
 \le d+k-n.
\]
Again \(m<k+d\), so \Cref{cor:finite-gluing} gives
\(T(k)\wedge U_{K/B'}(E)=0\).

At a height outside \(\supp(z)\), the finite representative of \(z\)
instead gives \(T(k)\wedge\Phi^Kz=0\), by
\Cref{prop:finite-local}. Thus every retained telescope kills
\(\Phi^K(z\otimes \mathcal A_B)\). This spectrum is \(L_m^f\)-local, so it is
zero. Forgetting a coherent action detects zero, which proves
\eqref{eq:outgoing-orthogonality}.

For \(J\not\sim B\), every summand of \(\Phi^B\mathcal A_J\) comes from a
proper normal subgroup \(J'\sim J\) of \(B\). Since
\(\ell(B/J',1)>0\), \Cref{cor:finite-gluing} at \(k=m\) kills that
summand after smashing with \(T(m)\). All smaller-height tests are killed
by the \(z\)-factor. The same detection argument proves
\eqref{eq:incoming-orthogonality}.
\end{proof}

\begin{theorem}[Relative specialization bound]\label{thm:specialization}
For a finite \(p\)-group \(G\), a subgroup \(H\le G\), and
\(m,n\ge0\),
\[
 (H,m)\in\overline{\{(G,n)\}}
 \quad\Longrightarrow\quad
 m\ge n+\ell(G,H).
\]
The closure is taken in the tensor-triangular spectrum of finite genuine
\(p\)-local \(G\)-spectra.
\end{theorem}

\begin{proof}
Fix \(G,n\) and put \(y=(G,n)\). Any specialization \((H,m)\) of
\(y\) has \(m\ge n\). Indeed, if \(m<n\), inflate a finite spectrum
of exact type \(m+1\). It belongs to \(\mathcal P_m(H)\) and not to
\(\mathcal P_n(G)\), contradicting \eqref{eq:prime-closure}.

Suppose that the theorem fails. Choose the least violating height \(m\)
and a violating subgroup \(B\) at that height. Thus
\[
 (B,m)\in\overline{\{y\}},\qquad m-n<\ell(G,B),
\]
and \(B\ne G\). By \Cref{prop:finite-local}, the spectrum
\(X_m=\Spc((\mathcal C_m^G)^c)\) is a finite subspace containing
\(y\) and \((B,m)\). Let \(Z=\overline{\{y\}}\cap X_m\), the
closure of \(y\) in that subspace.

We construct one compact object with support exactly \(Z\), keeping a
finite genuine representative. For every \(q\in X_m\setminus Z\),
reverse prime inclusion gives \(q\not\subseteq y\). Choose a finite
genuine spectrum \(a_q\in q\setminus y\). Its support contains \(y\),
hence all of \(Z\), and excludes \(q\). The finite tensor product
\[
 a=\bigotimes_{q\in X_m\setminus Z}a_q,\qquad z=e\otimes a
\]
(with the unit as empty product) satisfies \(\supp(z)=Z\).
It is compact and dualizable, and its displayed finite representative
allows us to apply \Cref{lem:layer-orthogonality}.

Minimality of \(m\) shows that every \((K,k)\in Z\) with \(k<m\)
satisfies \(\ell(G,K)\le k-n\). Therefore
\Cref{lem:layer-orthogonality,lem:layer-splitting} give
\[
 z\simeq X\oplus Q,\qquad
 Q=F_B\phi_Bz,\qquad X\otimes Q=0,
\]
with both summands compact, \(\phi_BX=0\), and
\(\phi_KQ=0\) for \(K\not\sim B\). Their compact supports are
disjoint closed subsets whose union is \(Z\). Both are nonempty:
\((B,m)\in\supp(Q)\), since \(\phi_BQ\simeq\phi_Bz\), and
\(y\in\supp(X)\), since \(B\ne G\) and \(\phi_GQ=0\).
This contradicts irreducibility of the closure of a point.
\end{proof}

\begin{proof}[Proof of \Cref{thm:main}]
For an integer \(r\ge0\), the universal implication defining
\(r_n(G,H)\) is precisely
\[
 \mathcal P_{n+r}^G(H)\subseteq\mathcal P_n^G(G).
\]
By \eqref{eq:prime-closure} and \Cref{thm:specialization}, this inclusion
forces \(r\ge\ell(G,H)\). Equivalently, for every smaller \(r\), its
failure supplies a finite genuine spectrum in the first prime and outside
the second, witnessing failure of the proposed implication.
\Cref{prop:cyclic-upper} supplies the opposite inequality. Thus
\(r_n(G,H)=\ell(G,H)\) for all the stated groups, subgroups, and heights.
The rational height \(n=0\) has been retained throughout, and when
\(H=G\) both sides are zero.
\end{proof}

\subsection{The general finite-group consequence}

For a finite group \(J\), let \(O^p(J)\) denote its smallest normal
subgroup of \(p\)-power index. We continue to use tested-height
indices and work in the \(p\)-local genuine category.

\begin{corollary}\label{cor:general-groups}
Let \(G\) be any finite group, \(p\) a prime, \(H,K\le G\), and
\(m,n\ge0\). Then
\[
 \mathcal P_m^G(K)\subseteq\mathcal P_n^G(H)
\]
if and only if there is a conjugate \(K^g=gKg^{-1}\le H\) such that
\[
 N=O^p(H)=O^p(K^g),\qquad
 m\ge n+\ell(H/N,K^g/N).
\]
In particular, the numerical test is applied to a qualifying
conjugate in the finite \(p\)-group \(H/N\).
\end{corollary}
\begin{proof}
Balmer--Sanders \cite[Proposition~6.11]{BS} reduce the displayed
inclusion to the existence of such a conjugate and the inclusion
\(\mathcal P_m^{H/N}(K^g/N)\subseteq
\mathcal P_n^{H/N}(H/N)\). Their prime indices are \(m+1,n+1\).
The common \(p\)-residual condition is equivalent to
\(p\)-subnormality of \(K^g\) in \(H\), and \(H/N\) is a
\(p\)-group. The relative specialization bound and the cyclic upper
bound give precisely the stated inequality. When \(m<n\), an
inflated finite spectrum of type \(m+1\) rules out the reduced
inclusion, as in the proof of \Cref{thm:specialization}.
For the upper direction at \(m>n+\ell(H/N,K^g/N)\), use also
the nesting of the Morava-acyclic ideals of finite spectra.
\end{proof}
